\section{Introduction}
\label{sec:introduction}
%% ====================================================================

Machine learning interatomic potentials
(MLIPs)~\cite{behler2007hdnnp,bartok2010gap} are models trained to predict the
energy and forces between atoms in a material, enabling simulations at near
quantum-mechanical accuracy but at a fraction of the computational cost of
\emph{ab initio} methods such as density functional theory
(DFT)~\cite{kohn1965dft}. Over the past decade, the field has produced a
growing ecosystem of MLIP algorithms---including
HDNNP~\cite{behler2007hdnnp}, GAP~\cite{bartok2010gap},
ACE~\cite{drautz2019ace}, MACE~\cite{batatia2022mace},
NequIP~\cite{batzner2022nequip}, M3GNet~\cite{chen2022m3gnet}, and
MTP~\cite{shapeev2016mtp}---each with distinct architectures, hyperparameters,
and application domains.

Researchers routinely face three interrelated challenges: selecting the right
algorithm for a given material and property, constructing suitable training
datasets, and optimizing
hyperparameters~\cite{grabowski2024metalearnerc}. These challenges are
compounded by a \emph{knowledge management problem}: the metadata connecting
MLIP studies---which algorithm was used, with which training set, for which
material system, at what accuracy---is scattered across published papers,
supplementary materials, ad-hoc scripts, and library-specific file formats.

Existing ontologies address parts of this landscape but not the specific needs
of MLIP research. ML ontologies (ML-Schema~\cite{publio2018mlschema},
DMOP~\cite{keet2015dmop}, MEX~\cite{esteves2015mex},
CMO~\cite{foltin2024cmo}) model general ML workflows but not
physics-specific hyperparameters or training set construction strategies.
Materials science ontologies (EMMO~\cite{horsch2020emmo},
MDO~\cite{li2024mdo}, PMDco~\cite{bayerlein2024pmdco},
CMSO/ASMO~\cite{menon2024atomrdf}) model the physical domain but not the
specifics of MLIP training. No existing ontology captures what is needed to
describe the provenance and composition of training datasets, the DFT settings
under which reference calculations were performed, and the algorithm-specific
hyperparameters that govern model accuracy.

In this paper, we present the \emph{\MLIPsOntology{}}, an ontology
expressed in the Web Ontology Language
(OWL\,2~DL)~\cite{owl2overview} that fills this gap. The ontology is organized into three modules:
%
\begin{itemize}
\item \emph{Algorithm}: captures MLIP algorithm families, their
  hyperparameters, implementations, and supported simulation types.
\item \emph{Training Data}: describes training datasets, their DFT
  provenance, material systems, and covered properties.
\item \emph{Benchmark}: links published evaluation studies to trained
  models, materials, and accuracy metrics.
\end{itemize}
%
The ontology is designed in alignment with ML-Schema, MDO, CMSO/ASMO, and
PROV-O, reusing their classes where appropriate and extending them with
MLIP-specific concepts. It includes 27 formal OWL axioms that enforce data
completeness, and SHACL shapes that validate data at the instance level.
The ontology is developed within the META-LEARN project, an ERC Advanced
Grant effort to create a meta-learning framework for MLIP
generation~\cite{grabowski2024metalearnerc}.

The remainder of this paper is organized as follows.
Section~\ref{sec:related-work} surveys related ontologies.
Section~\ref{sec:requirements} describes the requirements gathered from
domain experts.
Section~\ref{sec:ontology} presents the \MLIPsOntology{}, introduced through
a running example based on Moment Tensor Potentials.
Section~\ref{sec:evaluation} evaluates the ontology.
Section~\ref{sec:sustainability} discusses sustainability.
Section~\ref{sec:conclusion} concludes.

%% ====================================================================
